\section{A cap event produces a nonzero limit}
\label{sec:nonconcentration}
\label{nc:section}

The weak limit in Proposition~\ref{prop:blowup-limit} is obtained on an
open time interval.  A supremum attained by an approximation could in
principle disappear through concentration, or through escape from that
interval.  The purpose of this section is to exclude both possibilities.
The proof develops the nonconcentration mechanism of
\cite[Section~8]{OpenAIPeriodic}: two successive gains exclude an
axial limiting velocity, whereas five gains regularize the full
spacetime--velocity variables away from the axis. We give both
arguments here, including the localization needed for the rotated
period lattice. All constants may depend on the fixed jet parameters,
but none depends on the approximation index.

\begin{proposition}[Nonvanishing of every cap limit]
\label{prop:nonzero-limit}
Consider any sequence of cap solutions with horizons \(T_n\to0\),
approximation indices \(N_n\to\infty\), and
\(s_{N_n}/T_n\to0\).  Make the rescaling
\eqref{eq:rescaling}, and pass to any subsequence supplied by
Proposition~\ref{prop:blowup-limit}.  Its limit satisfies
\[
 g\not\equiv0\qquad\text{on }(0,U)\times\R^3_x\times\R^3_v.
\]
This conclusion includes \(l=1\), \(U=1\), and \(U=\infty\).
\end{proposition}

We first describe the positive kernels used in the proof.  Subsequently
we establish the two different mechanisms needed near and away from the
velocity axis.

\subsection{A positive majorizing equation and its kernels}

For this section relabel the sequence by \(n\), set \(U_n=T_n/r_n\),
and write
\[
 r_n=t_{*,n}+s_{N_n},\quad
 \ell_n=\frac{s_{N_n}}{r_n},\quad
 \tau_n=\frac{t_{*,n}}{r_n}=1-\ell_n,\quad
 h_n=h_{r_n},\quad C_n=C_{r_n},\quad f_n=\frac{h_n}{H}.
\]
Then \(\ell_n\to l\). In particular \(0\le\ell_n\le1\),
\(0\le\tau_n\le1\). The cap gives, on \(0\le t\le U_n\),
\begin{equation}
\label{nc:cap-normalized}
 0\le f_n(t,x,v)\le t+\ell_n,\qquad
 \left\|f_n(\tau_n)\right\|_\infty=1.
\end{equation}
Essential supremum norms and pointwise formulas involving \(H\) always
ignore the velocity axis, a null set.  The mild representatives of the
finite approximations will be used at the cap times.
Choose Borel representatives when composing densities with the maps
below, enforcing their asserted bounds on exceptional null sets as
well.  Characteristic identities and endpoint choices are still made
only on their common full-measure sets.

\begin{lemma}[Core kernels on the dilated scale]
\label{nc:core-kernels}
There are nonnegative velocity measures \(A_n(t,x,\dd p)\) and numbers
\(\varepsilon_n\to0\) with the following properties.
After discarding finitely many terms,
the measures \(A_n\) are defined from free transport on the whole interval
\(0\le t\le1\), independently of the cap interval \([0,U_n]\).
The bounds on \(A_n\) below hold on that whole interval; comparison with
\(h_n\) is used only while the cap bound holds.

\begin{enumerate}[label=\textup{(\roman*)}]
\item On \(0<t<\min(1,U_n)\), along almost every free characteristic,
\begin{equation}
\label{nc:positive-equation}
 \D h_n\le 2P(A_n,h_n)+\varepsilon_n H .
\end{equation}
At each point \(A_n\) contains at most one main free core.  Its dilated
scale \(\varrho=s/r_n\) satisfies \(|y|\asymp\varrho\), with coordinates measured
from its periodic copy, and
\[
 \frac{2P(A_n,H)}{H}\le\frac{C}{t+\ell_n}.
\]

\item Before spatial masks are imposed, a core of scale
\(\varrho=s/r_n\) and sign \(\sigma\) is the pushforward of
\begin{equation}
\label{nc:core-parameters}
 \frac{K S\bigl(\gamma(z-tp_3)/\varrho\bigr)}
      {\varrho+t\alpha\gamma S\bigl(\gamma(z-tp_3)/\varrho\bigr)}
 J(w)\dd w ,
 \qquad
 \bar p=d_s\bar w,\quad
 p_3=\sigma+\alpha\tanh\bigl(\gamma(z-tp_3)/\varrho\bigr)+d_sw_3 .
\end{equation}
The actual masks are measurable factors between zero and one.
The density in \eqref{nc:core-parameters} is at most
\(C(t+\varrho)^{-1}J(w)\).

\item For \(t\ge a>0\), the normalized gain measure
\begin{equation}
\label{nc:normalized-kernel}
 \mathcal K_n(t,x,v;\dd p\,\dd\xi)
 =
 4A_n(t,x,\dd p)\,
 \delta\bigl((v-p)\cdot(\xi-v)\bigr)
 \frac{H(\xi)}{H(v)}\dd\xi
\end{equation}
has total mass at most \(C/a\).  Uniformly in its base point, the
contributions from large \(|p|+|\xi-v|\), or from large retained
Gaussian parameters in a product-sheet majorant, tend to zero.
Its total mass tends uniformly to zero as \(|v|\to\infty\).
\end{enumerate}
\end{lemma}

\begin{proof}
Use the free-core measure from Lemma~\ref{dil:free-core-comparison}.
More explicitly, if the main retained label at \((r_nt,r_nx)\) is
\(i\), set
\[
 A_n(t,x,\dd p)=r_n f_i(r_nt,r_nx,p)\dd p,
 \qquad f_i(t,x,p)=c_i^0(x-tp,p),
\]
and otherwise set \(A_n=0\). The disjointness in
Lemma~\ref{cold:geometry} holds on the whole physical interval
\([0,r_n]\) for all sufficiently large \(n\), and no color is
retired there. Free transport therefore defines \(A_n\) on
\([0,1]\), regardless of \(U_n\). Equation~\eqref{dil:positive-comparison}
gives \eqref{nc:positive-equation}. Its error already includes
the bath, the source divided by \(m_{T_n}=T_n^p\), the rescaled
hot self-gain, and the difference between actual and free cold cores.

The enlargement of a main transverse annulus remains between its
original radial boundaries. Thus, relative to its period point,
\(\varrho<|y|<e^\eta\varrho\), where \(\varrho=s/r_n\). The change of velocity
variables in \eqref{cold:core-formula}, multiplied by \(r_n\), is
exactly \eqref{nc:core-parameters}; equivalently, this is the
parameter formula in \eqref{dil:free-core-integral} before velocity
testing. For \(0\le S\le1\),
\[
 \frac{KS}{\varrho+t\alpha\gamma S}
 \le \frac{C}{t+\varrho}.
\]
Indeed, multiplying by \(t+\varrho\) bounds the \(\varrho\)-term by \(K\)
and the \(t\)-term by \(K/(\alpha\gamma)\). The smallest scale is
\(\ell_n\), so the amplitude is at most \(C/(t+\ell_n)\).

The implicit axial map depends only on \(w_3\), and its monotonicity
gives \(|p_3|\le1+\alpha+d_s|w_3|\). The smooth compactly supported
\(J\) has a fixed product Gaussian majorant. Lemma~\ref{ker:parameter-sheets}
therefore dominates each core by a sheet, with the preceding amplitude
and a fixed axial moment bound. The gain bound in
Lemma~\ref{lem:sheet-bound} proves (i). The factor \(4\) in
\eqref{nc:normalized-kernel} is the product of the linear gain
coefficient \(2\) and the coefficient \(2\) in the collision-plane
formula of Lemma~\ref{lem:carleman}. Lemma~\ref{lem:kernel-tightness}
now gives (iii); Lemma~\ref{ker:parameter-sheets} also controls
Gaussian parameter tails before the compressed axial change of
variables. For the exact free cores the parameters already lie in
\(\supp J\). All these estimates allow measurable masks bounded by
one, which can only decrease the nonnegative integrals.
\end{proof}

For a fixed \(a>0\), define the backward Volterra operator
\begin{equation}
\label{nc:volterra}
 (\mathcal V_n F)(t,x,v)
 =
 \int_a^t\int
 F\bigl(s,x-(t-s)v,\xi\bigr)
 \mathcal K_n\bigl(s,x-(t-s)v,v;\dd p\,\dd\xi\bigr)\dd s .
\end{equation}
\emph{Representatives in the finite-temperature iteration.}
For each finite \(n\), the core measure has a velocity density,
\(A_n(s,x,\dd p)=a_n(s,x,p)\dd p\).
Integrating first in \(p\) in \eqref{nc:normalized-kernel} shows that,
for \(\xi\ne v\), its input marginal is
\[
\mathcal K_n(s,x,v;\R^3,\dd\xi)
=\frac{4H(\xi)}{|\xi-v|H(v)}
  \left[\int_{(v-p)\cdot(\xi-v)=0}
       a_n(s,x,p)\dd A_p\right]\dd\xi.
\]
There is no atom at \(\xi=v\): each plane-area measure gives that
point zero mass, and \(p=v\) is null for the cold density.
For a null set in \((s,X,\xi)\), integrate also in the endpoint
variables and use the measure-preserving change
\(X=x-(t-s)v\). Tonelli's theorem then shows that altering the input
on that set changes \(\mathcal V_nF(t)\) only on a null set of
endpoints, at each fixed terminal time. Finite repetition preserves
this property. A countable sequence of lower times \(a\downarrow0\)
suffices below; intersecting the good endpoint sets for these times,
the approximation indices, and the finitely many iterations gives
the common full-measure sets used at the cap events.

Positivity and Lemma~\ref{nc:core-kernels} give, for a bounded
nonnegative \(F\),
\begin{equation}
\label{nc:iteration-norm}
 \|\mathcal V_n^jF(t)\|_\infty
 \le \|F\|_\infty\frac{\bigl(C\log(t/a)\bigr)^j}{j!},
 \qquad a\le t\le1.
\end{equation}
Indeed, the ordered \(j\)-time integral of
\(\prod_{i=1}^j C/t_i\) equals
\((C\log(t/a))^j/j!\).  In particular every fixed number of further
gains has bounded mass on this slab.

\subsection{Reduction to a finite gain integral}

We begin the proof of Proposition~\ref{prop:nonzero-limit}. The case
\(l>0\) is settled directly. For \(l=0\), the cap will force a
positive contribution from five successive gains; the remaining
subsections show that this contribution vanishes if the weak limit is zero.


Pass to a subsequence on which the seed parameters converge,
\(\lambda_n\to\lambda_\infty\).  If \(l>0\) and
\(\lambda_\infty>0\), Proposition~\ref{prop:blowup-limit}
gives the initial hot trace
\[
 g(0,x,v)=\lambda_\infty l\,\chi(x/l,v).
\]
This trace is nonzero.  More explicitly, the limiting loss
frequency obeys
\[
 \nu_B(t,x,v)\le
       \frac{C(1+|v|)}{t+l}.
\]
On any bounded velocity set and any finite interval beginning
at zero its integral along a flight is finite.  The positive
mild formula for \eqref{eq:limit-equation} therefore retains
the initial trace multiplied by a strictly positive survival
factor.  A bounded positive portion of the smooth nonzero bump
gives positive interior mass.

If \(l>0\) and \(\lambda_\infty=0\), use
Lemma~\ref{nc:core-kernels} from time zero.  The initial ratio
has norm at most
\(\lambda_n\ell_n\|\chi/H\|_\infty\), independently of the
spatial scaling of the bump.  On \(0\le t\le\tau_n\le1\),
eventually \(t+\ell_n\ge l/2\).  The positive integral
inequality and the elementary scalar integrating-factor estimate
give
\[
 \sup_{0\le t\le\tau_n}\|f_n(t)\|_\infty
 \le
 \bigl(\lambda_n\ell_n\|\chi/H\|_\infty+\varepsilon_n\bigr)
 e^{2C/l}\longrightarrow0.
\]
This contradicts \eqref{nc:cap-normalized}.  Both alternatives
also apply when \(l=1\), even if \(\tau_n\to0\) or some hit
occurs at time zero.

It remains to consider \(l=0\), so \(\tau_n\to1\).
Suppose, to obtain a contradiction, that \(g=0\) throughout
the open interval \((0,U)\).  Nonnegativity upgrades this zero
weak limit to local \(L^1\) convergence: if \(K\) is a compact
subset of the open spacetime--velocity region, take a
nonnegative smooth compact test at least one on \(K\).  Its
pairing with \(h_n\) tends to zero, hence
\(\int_K h_n\to0\).  Away from the velocity axis \(H\) has a
positive lower bound on compact velocity sets.  Consequently
\begin{equation}
\label{nc:zero-in-measure}
 f_n\longrightarrow0
 \quad\text{in measure on compact subsets of }
 (0,U)\times\R^3_x\times\{\bar v\ne0\}.
\end{equation}

Fix \(0<a<1/4\), so that \(a<\tau_n\) for large \(n\).
The positive mild inequality \eqref{nc:positive-equation}
reads
\[
 f_n\le d_n+\mathcal V_nf_n+e_n,
\]
where \(d_n(t,x,v)=
 f_n(a,x-(t-a)v,v)\), \(\|d_n\|_\infty\le a+\ell_n\),
and \(0\le e_n\le\varepsilon_n\) on this unit-length interval.
Iterate the gain term five times.
Using positivity and \eqref{nc:iteration-norm} gives
\begin{equation}
\label{nc:five-residual}
 f_n(t,x,v)
 \le
 (a+\ell_n+\varepsilon_n)
 \sum_{j=0}^4\frac{(C\log(1/a))^j}{j!}
 +(\mathcal V_n^5 f_n)(t,x,v),
 \qquad a\le t\le\tau_n .
\end{equation}
Every formula is used on a common full-measure set of
characteristics for the finite approximation.


The term outside \(\mathcal V_n^5f_n\) becomes arbitrarily small when
first \(n\to\infty\) and then \(a\downarrow0\). We therefore need to
show that the five-fold term tends to zero at endpoints carrying a fixed
positive fraction of the cap. We first control those endpoints, then
separate axial and nonaxial limiting velocities.

\subsection{Compactness of the relevant endpoints}

\begin{lemma}[Endpoint and path truncations]
\label{nc:endpoint-tightness}
Fix \(a>0\), and let \(t_n\in[a,1]\).  In a five-fold integral
\((\mathcal V_n^5F_n)(t_n,x_n,v_n)\), assume \(0\le F_n\le2\).
For every prescribed error, one may restrict all
\(|p_i|+|v_i-v_{i-1}|\) and sheet parameters to fixed bounded sets at
that error, uniformly in \(n,t_n,x_n,v_n\).  Moreover:
\begin{enumerate}[label=\textup{(\roman*)}]
\item if \(|v_n|\to\infty\), the whole integral tends to zero;
\item after retaining bounded endpoint velocities, an endpoint
escaping every bounded neighborhood of every dilated periodic origin
also has a vanishing integral.
\end{enumerate}
Consequently any sequence along which the five-fold integral has a
positive lower bound has, after periodic translations and
subselection, convergent endpoint positions and velocities.
\end{lemma}

\begin{proof}
At any one gain, a removed normalized tail has uniformly small mass
by Lemma~\ref{nc:core-kernels}. Bound the other four gains and the
input by \eqref{nc:iteration-norm}, and sum over the five possible
locations of a removed tail. This proves the first assertion.
The same argument at the first gain, using its uniform decay for
large outgoing velocities, proves (i).

Retain bounded endpoint velocities and bounded increments. Then all
velocities are bounded by a fixed \(V\), and every collision position
is within distance \(V\) of the endpoint, since the total free-flight
time is at most one. Suppose
\[
 d_n:=\dist(x_n,r_n^{-1}\Lambda)\longrightarrow\infty.
\]
We show that the first kernel becomes uniformly small on these paths.
A main core of dilated scale \(\varrho\ge M\) has normalized mass at
most \(C/M\), by \eqref{nc:core-parameters} and the sheet bound.
Here \(M\) can be chosen arbitrarily large, independently of \(n\).
For a main core with \(\varrho\le M\), let \(q/r_n\), \(q\in\Lambda\),
be its period point. Write \((y,z)\) for the first collision position
relative to this point. Its transverse coordinate obeys
\(|y|\le e^\eta M\), while its distance to this point is at least
\(d_n-V\). Consequently \(|z|\to\infty\), uniformly over all
encountered cores of these scales. For bounded \(p_3,t\), their
parameter amplitudes satisfy
\begin{equation}
\label{nc:axial-escape}
 \frac{K S(\gamma(z-tp_3)/\varrho)}
      {\varrho+t\alpha\gamma S(\gamma(z-tp_3)/\varrho)}
 \le \frac K\varrho S(\gamma(z-tp_3)/\varrho)\longrightarrow0
\end{equation}
uniformly for \(0<\varrho\le M\). In fact, when
\(|z-tp_3|\ge Z\), the right side is at most
\(4K\varrho^{-1}e^{-2\gamma Z/\varrho}\), whose supremum tends to zero
as \(Z\to\infty\). Applying the sheet bound with this vanishing
amplitude proves the assertion for \(\varrho\le M\). The remaining
gains have bounded mass. First let \(n\to\infty\), then
\(M\to\infty\), and finally remove the initial tail error.
This proves (ii).

It follows that a sequence with a positive lower bound for its
five-fold integral has bounded velocities and bounded distance from
\(r_n^{-1}\Lambda\), after subselection. Choose period points
\(q_n/r_n\) realizing a bounded distance, and translate the endpoints
by these periods. The resulting positions have a convergent
subsequence. Such translations preserve both the kernels and the
periodic hot inputs used below.

We record also the geometry along the translated bounded paths.
Every physical main enlargement for times \(r_nt\le r_n\) lies
in a fixed ball of radius less than \(1/2\) about its period point,
for sufficiently small \(s_0,r_n\). Distinct points of \(\Lambda\)
are at least distance one apart. A path in a bounded dilated set
lies at physical distance \(O(r_n)\) from the chosen origin, so
for all large \(n\) it can encounter only that copy. Its encountered
scales satisfy \(s/r_n\le C\), by the transverse annular restriction.
Thus the largest physical width along these paths is
\(O(r_n^{10})\), which tends to zero.
\end{proof}

\subsection{Kernel estimates near the velocity axis}

The preceding localization reduces the problem to bounded endpoint
velocities. To distinguish axial and nonaxial limits, we now refine
the normalized sheet bounds of Section~\ref{sec:kernel}. The estimates
below, from \cite[Section~3]{OpenAIPeriodic}, allow transverse
concentration of the hot input but prevent concentration in its axial
coordinate. We recall that $\mathcal S_L$ is the class of probability
mixtures of sheets $\mathcal G_d\otimes\rho$, where
$0\leq d\leq1/8$ and $\int e^{8a^2}\rho(\dd a)\leq L$, and that
$\mathcal K_c$ denotes the normalized measure in
\eqref{ker:normalized-kernel}.

\begin{lemma}[Kernel estimates near the axis]
\label{lem:axis-kernel-estimates}
Fix $1\leq L<\infty$. The following conclusions hold uniformly for
$c\in\mathcal S_L$ and for outgoing velocities off the axis.
\begin{enumerate}[label=\textup{(\roman*)}]
\item For every fixed $M,\delta>0$,
\begin{align}
 &\lim_{\varepsilon\downarrow0}
 \sup_{c\in\mathcal S_L}
 \esssup_{\substack{|v|\leq M\\0<|\bar v|\leq\varepsilon}}
 \mathcal K_c(v;|\bar\xi|\geq\delta)=0,
 \label{ker:axis-input}\\
 &\lim_{\varepsilon\downarrow0}
 \sup_{c\in\mathcal S_L}
 \esssup_{\substack{|v|\leq M\\0<|\bar v|\leq\varepsilon}}
 \mathcal K_c(v;|v_3-p_3|\geq\delta)=0.
 \label{ker:axis-mismatch}
\end{align}
\item The axial input marginals are uniformly non-atomic. More
precisely, for each $R\geq1$ and each interval $I$ of length
$0<\ell\leq1$,
\begin{equation}\label{ker:axial-interval}
 \mathcal K_c(v;|p|+|\xi-v|\leq R,\ \xi_3\in I)
 \leq C_{q,m,L,R}\bigl(\ell+\ell^{2-q}\bigr).
\end{equation}
In combination with \eqref{ker:uniform-tail}, this gives
\begin{equation}\label{ker:axial-modulus}
 \lim_{\ell\downarrow0}\sup_{c\in\mathcal S_L}\esssup_v
 \sup_{|I|\leq\ell}\mathcal K_c(v;\xi_3\in I)=0.
\end{equation}
\item Let $\mathcal S_L(\varepsilon)$ denote the subclass whose
elementary transverse widths are at most $\varepsilon$. For every
$M<\infty$ and $r_0>0$,
\begin{equation}\label{ker:off-axis-small-input}
 \lim_{\delta\downarrow0}\limsup_{\varepsilon\downarrow0}
 \sup_{c\in\mathcal S_L(\varepsilon)}
 \esssup_{\substack{|v|\leq M\\|\bar v|\geq r_0}}
 \mathcal K_c(v;|\bar\xi|<\delta)=0.
\end{equation}
\end{enumerate}
All these conclusions hold under uniformly bounded sheet amplitudes.
\end{lemma}

\begin{proof}
By \eqref{ker:uniform-tail}, we may work on
$|p|+|\xi-v|\leq R$ and remove that restriction at the end of each
argument. On this region the exponential and polynomial factors in
\eqref{ker:energy-factor} are bounded by a constant depending on $R$.

For \eqref{ker:axis-input}, on $|\bar\xi|\geq\delta$ the factor
$W(\xi)$ is bounded by $1+\delta^{-q}$. The regular line estimate
and \eqref{ker:transverse-expectation} therefore give
\[
 C_{R,\delta}\frac{|\bar v|^{-1}}{1+|\bar v|^{-q}}
 \leq C_{R,\delta}|\bar v|^{q-1}
 \qquad(0<|\bar v|\leq1).
\]
For \eqref{ker:axis-mismatch}, integrate in $\xi_3$ first, using
$|v_3-p_3|^{-1}\leq\delta^{-1}$ and
\eqref{ker:disk-singularity}. The bound is
$C_{R,\delta}/W(v)\leq C_{R,\delta}|\bar v|^q$.
These are uniform in the axial law. For
\eqref{ker:axial-interval}, discard the bounded exponential and
polynomial factors and apply
\eqref{ker:geometric-interval}. This also proves
\eqref{ker:axial-modulus} after the tail removal.

It remains to prove \eqref{ker:off-axis-small-input}. Fix the
truncation radius $R$. The part $|\bar p|>r_0/4$ is uniformly
negligible as the transverse widths tend to zero. Indeed,
\[
 \ind_{|\bar p|>r_0/4}\mathcal G_d(\dd\bar p)
 \leq2e^{-r_0^2/(32d^2)}\mathcal G_{\sqrt2d}(\dd\bar p),
\]
and Lemma~\ref{lem:sheet-bound} applies to the wider Gaussian for
all sufficiently small $d$. Now retain $|\bar p|\leq r_0/4$.
Write $n=\bar v-\bar p$, $A=(n/|n|)\cdot\bar v$, and
$B=(v_3-p_3)/|n|$, as in Lemma~\ref{ker:geometric-slice}. Then
\[
 |n|\geq\tfrac34r_0,\qquad A\geq\tfrac35r_0.
\]
If $|v_3-p_3|\leq r_0^2/(8R)$, then for
$|\xi_3-v_3|\leq R$,
\[
 |B(\xi_3-v_3)|\leq\tfrac16r_0.
\]
The transverse collision line stays at distance greater than
$r_0/3$ from zero, and cannot meet $|\bar\xi|<\delta$ when
$\delta<r_0/3$. In the complementary case,
$|v_3-p_3|^{-1}\leq8R/r_0^2$. Integrating first in $\xi_3$ bounds
the small-input contribution by
\[
 C_{R,r_0}\int_{|\bar\xi|<\delta}
               (1+|\bar\xi|^{-q})\,\dd\bar\xi
 \leq C_{R,r_0}(\delta^2+\delta^{2-q}).
\]
First let the widths tend to zero, then $\delta$ tend to zero, and
finally remove the truncation. This proves the last assertion.
\end{proof}

These conclusions apply to the core kernels
\eqref{nc:normalized-kernel}: for $t\geq a>0$ the proof of
Lemma~\ref{nc:core-kernels} gives sheet domination with a fixed axial
moment bound and amplitude $C/a$. On the localized paths, the
transverse widths tend uniformly to zero, as shown at the end of
Lemma~\ref{nc:endpoint-tightness}. This is the additional hypothesis
needed for \eqref{ker:off-axis-small-input}.

\begin{remark}[A sharp axial endpoint]\label{ker:matching-axis}
A single normalized gain need not become small as the outgoing
velocity approaches the axis. Recall
$a_q=\int_{\R}(1+s^2)^{-q/2}\dd s$ from \eqref{ker:line-integral}. Take
$c=\delta_{(0,0,u)}$ and
$v_\varepsilon=(\varepsilon,0,u)$. The collision plane is
$\xi_1=\varepsilon$, and
\[
 \frac{P(c,H)(v_\varepsilon)}{H(v_\varepsilon)}
 =\frac{2}{\varepsilon H(v_\varepsilon)}
      \int_{\R^2}H(\varepsilon,y,z)\,\dd y\,\dd z.
\]
The regular part of $W(\varepsilon,y,z)$ contributes
$O(\varepsilon^{q-1})$. In its singular part substitute
$y=\varepsilon s$. Dominated convergence, with dominating function
$(1+s^2)^{-q/2}e^{-z^2}\langle z\rangle^{-m}$, gives
\[
 \lim_{\varepsilon\downarrow0}
 \frac{P(c,H)(v_\varepsilon)}{H(v_\varepsilon)}
 =2a_q e^{u^2}\langle u\rangle^m
        \int_{\R}e^{-z^2}\langle z\rangle^{-m}\,\dd z>0.
\]
The corresponding axial input marginal converges to the absolutely
continuous measure
\[
2a_q e^{u^2}\langle u\rangle^m
 e^{-z^2}\langle z\rangle^{-m}\,\dd z.
\]
Thus transverse concentration is compatible with the uniform axial
interval estimate. This is the reason for using two successive gains
in the axial argument that follows.
\end{remark}

\subsection{Two gains at an axial endpoint}

For \(\varrho>0\) and \(\sigma\in\{-1,1\}\), denote by
\(u_{\varrho,\sigma}(t,z)\) the solution of
\[
 u=\sigma+\alpha\tanh\bigl(\gamma(z-tu)/\varrho\bigr).
\]
This notation extends the profile independently of every annular
or stripe mask.

\begin{lemma}[Characteristic matching and small scales]
\label{nc:characteristic-matching}
On every compact set with \(t\ge a>0\), the following statements
hold uniformly.
\begin{enumerate}[label=\textup{(\roman*)}]
\item If \(a\le t_2\le t_1\), \(\Delta=t_1-t_2\), and
\(z_2=z_1-\Delta w\), then
\begin{equation}
\label{nc:matching-inequality}
 |w-u_{\varrho,\sigma}(t_1,z_1)|
 \le |w-u_{\varrho,\sigma}(t_2,z_2)|.
\end{equation}
\item As \(\varrho\downarrow0\), the distance of
\(u_{\varrho,\sigma}(t,z)\) to
\[
 \{z/t,\ \sigma-\alpha,\ \sigma+\alpha\}
\]
tends to zero uniformly.  For the perturbed axial velocity in
\eqref{nc:core-parameters},
\begin{equation}
\label{nc:thermal-speed-error}
 |p_3-u_{\varrho,\sigma}(t,z)|\le d_s|w_3|.
\end{equation}
\end{enumerate}
\end{lemma}

\begin{proof}
The profile satisfies
\[
 \partial_z u_{\varrho,\sigma}
 =\frac{(\alpha\gamma/\varrho)S}
        {1+(t\alpha\gamma/\varrho)S}\ge0,
 \qquad
 \partial_tu_{\varrho,\sigma}
       +u_{\varrho,\sigma}\partial_zu_{\varrho,\sigma}=0.
\]
For fixed \(t_1,t_2,z_1\), the function
\(F(w)=w-u_{\varrho,\sigma}(t_2,z_1-\Delta w)\) has derivative
\(F'(w)\ge1\).  Its root is
\(u_{\varrho,\sigma}(t_1,z_1)\): the defining implicit equation shows
that this velocity is unchanged along the straight segment between
the two points.  Integrating \(F'\ge1\) proves
\eqref{nc:matching-inequality}.

Fix \(0<\delta<\alpha\).  If \(u_{\varrho,\sigma}\) is within
\(\delta\) of an endpoint speed, there is nothing to prove.
Otherwise
\[
 \left|\frac{u_{\varrho,\sigma}-\sigma}{\alpha}\right|
 \le1-\frac{\delta}{\alpha},
 \qquad
 |z-tu_{\varrho,\sigma}|
 \le\frac{\varrho}{\gamma}
       \operatorname{arctanh}(1-\delta/\alpha).
\]
It follows that
\[
 |u_{\varrho,\sigma}-z/t|
 \le\frac{\varrho}{\gamma a}
       \operatorname{arctanh}(1-\delta/\alpha).
\]
First choose \(\delta\), then \(\varrho\) small.  This proves the
uniform assertion in (ii).
Finally, the function
\[
 p\longmapsto p-\sigma-\alpha\tanh\bigl(\gamma(z-tp)/\varrho\bigr)
\]
has derivative at least one.  Its values at \(u_{\varrho,\sigma}\)
and at the perturbed \(p_3\) are \(0\) and \(d_sw_3\).
This proves \eqref{nc:thermal-speed-error}.
\end{proof}

\begin{lemma}[Vanishing of two axial gains]
\label{nc:two-axial-gains}
Fix \(a>0\).  Let \(t_n\in[a,1]\), and suppose the endpoint
positions and velocities are bounded with \(\bar v_n\to0\).
Then
\[
 (\mathcal V_n^2 1)(t_n,x_n,v_n)\longrightarrow0.
\]
The conclusion continues to hold if all subsequent factors in a
longer finite gain integral are bounded uniformly.
\end{lemma}

\begin{proof}
Remove normalized parameter and increment tails at an arbitrarily
small cost, using Lemma~\ref{nc:core-kernels}.  On the retained
two-gain paths, all velocities and positions lie in fixed compact
sets, the times lie in \([a,1]\), and every active scale satisfies
\(0<\varrho\le M\).  In particular the largest possible physical width
is \(d_0(r_nM)^{10}\to0\).

Write \(t_1,x_1=(y_1,z_1)\) for the first collision point, and
\(\xi\) for its input velocity.  The next point has axial coordinate
\(z_2=z_1-(t_1-t_2)\xi_3\).  All estimates below are uniform when
\(t_1,x_1\) vary in their retained compact sets.  At the first
gain, Lemma~\ref{lem:axis-kernel-estimates} supplies a common modulus
\(\omega(\rho)\downarrow0\) for the mass of any axial interval of
length \(\rho\); the interval's center is arbitrary.  It also says
that for each fixed \(b>0\), the first-kernel mass on
\(\{|\bar\xi|>b\}\) tends to zero as \(\bar v_n\to0\).

Here is an order of choices that handles all scales at the second
point.  Fix an error \(\epsilon>0\), including the fixed factors
from the two time integrations and the bound \(C/a\).  Choose
\(\delta>0\) so that the sum of the first-kernel masses of the
\(\delta\)-neighborhoods of the five numbers
\begin{equation}
\label{nc:collapsed-speeds}
 z_1/t_1,\quad 1-\alpha,\quad1+\alpha,\quad
 -1-\alpha,\quad-1+\alpha
\end{equation}
is at most \(\epsilon\).  This is possible uniformly because
there are five intervals and their lengths are \(2\delta\).
Lemma~\ref{nc:characteristic-matching} then gives
\(\varrho_0>0\), chosen uniformly on the retained compact sets, so that
every profile with \(\varrho<\varrho_0\) is within \(\delta/4\) of its
corresponding three collapsed speeds.

The geometric mesh has a uniformly bounded number \(J_0\) of
scales in \([\varrho_0,M]\).  Choose \(b_0>0\) so small that the
first-kernel mass of the union of the \(b_0\)-neighborhoods of
\[
 \{u_{\varrho,\sigma}(t_1,z_1):
       \varrho\in[\varrho_0,M]\text{ in the current mesh},\
       \sigma=\pm1\}
\]
is at most \(\epsilon\).  For example it suffices that
\(2J_0\omega(2b_0)\le\epsilon\), after incorporating the fixed
amplitude bound into \(\omega\).  This list may depend on \(n\);
its cardinality bound does not.  Its centers may depend on
\(t_1,z_1\), which are held fixed when the first kernel is
integrated.  Thus neither dependence weakens the interval bound.

Let \(E_n(t_1,z_1)\) be the union of the neighborhoods of the five
collapsed speeds in \eqref{nc:collapsed-speeds} and the neighborhoods
of the finite-scale profile speeds just chosen.
Outside \(E_n\), every retained cold speed at the second point has
an axial mismatch bounded below by a fixed positive number for
all large \(n\).  For \(\varrho\ge\varrho_0\), this follows from
\eqref{nc:matching-inequality} and
\eqref{nc:thermal-speed-error}, with lower bound \(b_0/2\).
For a constant collapsed speed it follows immediately, with lower
bound \(\delta/2\).  For the remaining branch use the exact identity
\[
 \xi_3-\frac{z_2}{t_2}
 =\frac{t_1}{t_2}
       \left(\xi_3-\frac{z_1}{t_1}\right);
\]
the factor \(t_1/t_2\ge1\) preserves the lower bound.
The errors \(\delta/4\) in small-scale approximation and the
vanishing thermal errors can be absorbed.  We have therefore
obtained, on the complement of \(E_n\),
\[
 |\xi_3-p_3^{(2)}|
 \ge c_*:=\tfrac14\min(\delta,b_0)>0
\]
for every retained second-core parameter.

The second-kernel axial mismatch estimate in
Lemma~\ref{lem:axis-kernel-estimates} now tends uniformly to zero as
\(|\bar\xi|\to0\).  Choose \(b>0\) so small that this second
kernel has mass at most \(\epsilon\) whenever
\(|\bar\xi|\le b\) and the displayed mismatch holds.
For this fixed \(b\), the first-kernel mass on
\(|\bar\xi|>b\) tends to zero.  Its mass on \(E_n\) is at most
\(2\epsilon\), while the remaining first-kernel mass is bounded.
Combining the three regions, then removing the initial tail error,
gives a bound tending to zero as \(\epsilon\downarrow0\).
All choices preceded the final passage \(n\to\infty\), and were
uniform in the two times.  This proves the lemma.

In a longer integral, all factors following these two gains are
bounded by \eqref{nc:iteration-norm}; the same proof applies.
\end{proof}

\subsection{Finite-dimensional regularization away from the axis}

Two gains have now ruled out endpoints approaching the velocity axis.
At a nonaxial endpoint the kernel mass need not vanish. Instead, we
will show that five gains cannot concentrate their input on a set of
small spacetime--velocity volume. The next three lemmas make this
statement precise: they parameterize the collision planes, give an
abstract estimate for composing measurable functions with smooth maps,
and verify the required rank for five successive collisions.

\begin{lemma}[Redundant plane parameters]
\label{nc:plane-parameters}
Let \(n\ne0\), \(e=n/|n|\), and let \(\zeta\) be a bounded
nonnegative function supported in \([-1,1]\) with integral one.
For every nonnegative measurable \(F\),
\begin{equation}
\label{nc:plane-identity}
 \int_{n^\perp} F(k)\frac{\dd A(k)}{|n|}
 =
 \int_{\R^3}F(n\times L)|n|\zeta(L\cdot e)\dd L .
\end{equation}
When \(\delta\le|n|\le R\) and the plane integral is restricted
to \(|k|\le R_1\), the corresponding parameter measure
\(\ind_{\{|n\times L|\le R_1\}}|n|\zeta(L\cdot e)\dd L\)
is supported in
\[
 |L|\le\sqrt{(R_1/\delta)^2+1}
\]
and has density at most \(R\|\zeta\|_\infty\). This bound concerns
the parameter measure before multiplication by \(F(n\times L)\).
The identity can be applied successively when each normal is a
measurable function of parameters chosen at earlier gains.
\end{lemma}

\begin{proof}
Write \(L=L^\perp+\lambda e\).  The map
\(L^\perp\mapsto n\times L^\perp\) multiplies two-dimensional
area by \(|n|^2\).  Integration in \(\lambda\) contributes
\(\int\zeta=1\), proving \eqref{nc:plane-identity} and the support
and density bounds.  For successive gains, fix all previous
parameters before applying the identity.  The current normal is
then fixed; Fubini's theorem for nonnegative functions justifies
iteration.  No derivatives of the variable normals enter this
conditional change of variables.
\end{proof}

\begin{lemma}[Uniform pullback estimate]
\label{nc:uniform-pullback}
Let \(\mathcal O\subset\R^k\) be open, \(k\ge d\), and suppose
\(\Phi_n,\Phi:\mathcal O\to\R^d\) satisfy
\(\Phi_n\to\Phi\) in \(C^1_{\mathrm{loc}}\) and
\(\operatorname{rank}D\Phi=d\) almost everywhere.
Let \(K\Subset\mathcal O\) and \(B\subset\R^d\) be compact.
Let \(a_n\) be measurable, supported in \(K\), with
\(|a_n|\le A\).  Let \(F_n\) be measurable, supported in \(B\),
with \(|F_n|\le M\), and converging to zero in measure.  Then
\begin{equation}
\label{nc:pullback-conclusion}
 \int_{\mathcal O}|F_n(\Phi_n(q))a_n(q)|\dd q\longrightarrow0 .
\end{equation}
No convergence or continuity of the densities \(a_n\) is required.
\end{lemma}

\begin{proof}
Boundedness and convergence in measure on the finite-volume set
\(B\) imply \(\|F_n\|_{L^1(\R^d)}\to0\): for every
\(\eta>0\), bound the integral by
\(\eta|B|+M|\{|F_n|>\eta\}|\).

Fix \(\epsilon>0\).  The critical set in \(K\) has measure zero,
so choose an open neighborhood \(E\) of it with
\(|E\cap K|<\epsilon\).  The compact set \(K\setminus E\)
consists of regular points.  At such a point \(q_0\), choose \(d\)
input coordinates giving a nonzero minor of \(D\Phi(q_0)\).
Let \(\pi\) be projection onto the complementary \(k-d\)
coordinates, and define
\[
 G(q)=(\Phi(q),\pi q),\qquad
 G_n(q)=(\Phi_n(q),\pi q),\qquad A_0=DG(q_0).
\]
The matrix \(A_0\) is invertible.  Choose a convex ball
\(V\Subset\mathcal O\) about \(q_0\) so small that, for all
sufficiently large \(n\),
\[
 \sup_{q\in V}\|A_0^{-1}DG_n(q)-I\|<\tfrac12.
\]
Integration on the line segment between \(q,q'\in V\) gives
\[
 |A_0^{-1}(G_n(q)-G_n(q'))-(q-q')|
 \le\tfrac12|q-q'|.
\]
Thus \(G_n\) is injective on \(V\), and its Jacobian has a
positive lower bound \(c_V\) independent of \(n\).  The inverse
function theorem consequently makes it a diffeomorphism onto
its image.

Cover \(K\setminus E\) by finitely many such balls \(V_j\).
The change-of-variables formula gives
\[
 \begin{split}
 \int_{V_j\cap K}|F_n(\Phi_n(q))|\dd q
 &\le c_{V_j}^{-1}
       \int_{G_{n,j}(V_j\cap K)}|F_n(z)|\dd z\,\dd w\\
 &\le c_{V_j}^{-1}|\pi_j(V_j)|
       \|F_n\|_{L^1(\R^d)}.
 \end{split}
\]
When \(k=d\), the zero-dimensional factor has value one.
The integral on \(E\cap K\), including its density, is at most
\(AM\epsilon\). Multiplying the regular-chart estimates by \(A\)
and summing over the finite cover gives
\[
 \limsup_{n\to\infty}
 \int_{\mathcal O}|F_n(\Phi_n(q))a_n(q)|\dd q
 \le AM\epsilon.
\]
Let \(\epsilon\downarrow0\).
This proves \eqref{nc:pullback-conclusion}. The argument used the
densities only through \(|a_n|\le A\), so arbitrary measurable masks
and restrictions may be included in them without any convergence
assumption.
\end{proof}

\begin{lemma}[Five-gain rank and robust pullback]
\label{lem:five-gain-rank}
Fix a nonaxial endpoint velocity \(v_0\), an endpoint
\((t_0,x_0)\) with \(t_0>0\), five positive scales
\(\varrho_i\), and five signs \(\sigma_i\).
Given positive times \(t_i\) and \(L_i\in\R^3\), define recursively
\begin{equation}
\label{nc:rank-map}
 \begin{split}
 x_i&=x_{i-1}-(t_{i-1}-t_i)v_{i-1},\\
 p_i&=u_{\varrho_i,\sigma_i}(t_i,z_i)e_3,\qquad
 n_i=v_{i-1}-p_i,\\
 k_i&=n_i\times L_i,\qquad v_i=v_{i-1}+k_i ,
 \qquad i=1,\ldots,5 .
 \end{split}
\end{equation}
The map
\[
 \Phi:(t_1,\ldots,t_5,L_1,\ldots,L_5)
       \longmapsto(t_5,x_5,v_5)
\]
is real analytic on the connected open set
\((0,\infty)^5\times\R^{15}\), and has rank seven almost everywhere.
In particular it is a submersion almost everywhere on the ordered
region \(0<t_5<\cdots<t_1<t_0\).

For convergent endpoint sequences with a nonaxial limiting velocity,
convergent positive scale sequences, and the finite-temperature
core formulas \eqref{nc:core-parameters}, the analogous maps converge
in \(C^1\) on compact parameter sets to \(\Phi\).
Bounded measurable parameter densities supported in a common
compact set therefore satisfy
Lemma~\ref{nc:uniform-pullback}, including after any measurable
spatial, velocity, or time restrictions.
\end{lemma}

\begin{proof}
For \(t>0\), the derivative with respect to \(u\) in the defining
implicit equation is
\[
 1+\frac{t\alpha\gamma}{\varrho}
 S\bigl(\gamma(z-tu)/\varrho\bigr)>0.
\]
The implicit analytic function theorem gives an analytic profile
for every \(z\in\R\).  Uniqueness, which also follows from this
strict monotonicity, identifies the local profiles on overlaps.
The recursive formulas \eqref{nc:rank-map} are consequently
analytic.  They make sense even without ordering the positive
times: a negative formal free-flight duration merely changes
the point at which the next analytic profile is evaluated.
This supplies the connected open domain stated in the lemma.

We exhibit one rank-seven point.  Fix \(t_5=\tau\) with
\(0<\tau<t_0\), and set all five collision times equal to
\(\tau\).  After a rotation about the velocity axis, write
\(v_0=(R,0,z)\), \(R>0\).  All collision positions now coincide
at \(x_0-(t_0-\tau)v_0\), independently of the angular parameters.
Thus write the five axial velocities simply as \(p_i=a_i e_3\);
their values \(a_i\) may be arbitrary.

Choose
\[
 k_1=(0,1,0),\qquad k_2=(1,-R,0).
\]
They satisfy \(k_1\perp n_1\), \(k_2\perp n_2\), and span the
transverse plane.  Put
\[
 w=\bar v_2=(R+1,1-R)\ne0,\qquad b_3=v_{2,3}-a_3,
\]
and choose a sufficiently small nonzero number \(\epsilon\) and
\[
 k_3=\epsilon\left(-\frac{b_3w}{|w|^2},1\right).
\]
Then \(k_3\perp n_3\), its axial component is nonzero, and
\(\bar v_3\ne0\).  Hence \(k_1,k_2,k_3\) are linearly independent.
Writing \(\bar v_3=(X,Y)\), take
\[
 k_4=(-Y,X,0),\qquad k_5=0.
\]
The fourth increment is perpendicular to \(n_4\), and
\[
 \det(\bar n_4,\bar n_5)=X^2+Y^2>0.
\]
Thus \(n_4,n_5\) are nonparallel.  All the normals in this
construction are nonzero, and every chosen increment can be
represented by \(L_i=-n_i\times k_i/|n_i|^2\).  Choose \(L_5=0\).

At the equal-time configuration,
\[
 x_5=x_0-v_0(t_0-t_1)
           -\sum_{i=1}^4v_i(t_i-t_{i+1}).
\]
All derivatives of \(x_5\) in angular directions vanish.
Keeping \(t_5\) fixed, its first three time columns are
\[
 \partial_{t_1}x_5=-k_1,\qquad
 \partial_{t_2}x_5=-k_2,\qquad
 \partial_{t_3}x_5=-k_3.
\]
Derivatives of the \(v_i\), including derivatives of the
position-dependent \(p_i\), are multiplied by zero free-flight
durations in this formula.  In the final velocity block,
\[
 \delta_{L_4}v_5=n_4\times\delta L_4,\qquad
 \delta_{L_5}v_5=n_5\times\delta L_5,
\]
because \(L_5=0\) and the \(p_i\) have zero angular position
derivatives at equal times.  The sum of the planes
\(n_4^\perp+n_5^\perp\) is \(\R^3\).  Choose three angular
columns spanning this sum.  Together with the three time columns
above, they give an invertible \(6\times6\) block for
\((x_5,v_5)\): its upper-right block is zero and both diagonal
blocks are invertible.  Appending the \(t_5\) column and target
coordinate gives a nonzero \(7\times7\) minor.

This minor is a nontrivial analytic function on the connected
domain.  Its zero set has Lebesgue measure zero.  For completeness,
at any zero of a nontrivial analytic function some derivative of
finite order is nonzero; otherwise its Taylor series vanishes on
a neighborhood, and analytic continuation through overlapping
balls in the connected domain makes the function identically
zero.  Choosing a derivative of minimal nonzero order puts the
point on a regular zero set of a derivative of one lower order.
There are countably many such derivatives, and each regular zero
set is covered by countably many smooth hypersurface charts.
Their union is null.  Rank failure is contained in the zero set
of our particular minor, so it is null as well.  Continuity also
shows directly that the witness persists at nearby strictly
ordered times.  It need not obey any spatial mask.

For the prelimit maps replace \(p_i\) by
\eqref{nc:core-parameters} and include \(w_i\in\R^3\) among
the parameters.  On compact sets the scales have positive lower
bounds, \(d_{r_n\varrho_i}\to0\), and all implicit derivatives have
the same positive lower bound one.  Differentiating the implicit
equations, or applying their implicit function formulas, gives
\(C^1\) convergence in the times, positions, and all bounded
\(L_i,w_i\).  The limit is \(\Phi\), independent of the fifteen
spectator \(w_i\)-coordinates.  Its rank is still seven almost
everywhere on the enlarged parameter domain, by Fubini.
Lemma~\ref{nc:uniform-pullback} now proves the final assertion.
\end{proof}

\begin{lemma}[Vanishing of five nonaxial gains]
\label{nc:five-nonaxial-gains}
Suppose \(\ell_n\to0\), the hit endpoints converge after periodic
translation, and their limiting velocity is nonaxial.  Suppose
also that \(0\le f_n\le2\) on \(a\le t\le\tau_n\) and
\(f_n\to0\) in measure on every compact subset of
\[
 (0,U)\times\R^3_x
       \times\{v\in\R^3:\bar v\ne0\}.
\]
Then
\[
 (\mathcal V_n^5f_n)(\tau_n,x_n,v_n)\longrightarrow0.
\]
\end{lemma}

\begin{proof}
Fix an arbitrarily small error tolerance.  Remove the tails of
Lemma~\ref{nc:endpoint-tightness}.  All path velocities and
positions are now in fixed bounded sets.  At every encountered
main core, \(s\le Cr_n\), so its transverse width tends uniformly
to zero.

Starting from a positive lower bound for the endpoint transverse
speed, apply the nonaxial small-input estimate in
Lemma~\ref{lem:axis-kernel-estimates} successively at gains
\(1,\ldots,5\).  Choose positive numbers \(b_1,\ldots,b_5\)
so that deleting \(|\bar v_i|<b_i\) costs less than the prescribed
error at step \(i\).  At each step the previous lower bound has
already been fixed; the later gain masses are bounded by
\eqref{nc:iteration-norm}.  Thus this is a finite successive
choice, followed by \(n\to\infty\).  The sum of the five deleted
contributions is arbitrarily small.  Write \(b>0\) for a lower
bound on all retained transverse speeds, including \(v_0\).
For large \(n\), every retained normal \(v_{i-1}-p_i\) has length
at least \(b/2\).

Next delete collision points with \(|y_i|<\rho\).
Conditional on all earlier parameters, \(y_i\) is an affine
function of \(t_i\) with derivative \(\bar v_{i-1}\).
Consequently
\[
 \left|\{t_i:|y_i|<\rho\}\right|\le\frac{2\rho}{b}.
\]
Integrate the current kernel using its mass bound \(C/a\),
and bound the later kernels by \eqref{nc:iteration-norm}.
Summing over the five steps bounds this deletion by \(C_a\rho/b\).
Choose \(\rho>0\) after the speed thresholds.  The retained
scales now belong to a compact subinterval
\([\varrho_-,\varrho_+]\subset(0,\infty)\).
Mesh convergence gives a fixed finite list of scales on this
interval, with positive limits, after an immaterial enlargement
of its endpoints.  There are consequently only finitely many
five-tuples of scales and signs to consider.

The tested time \(t_5\) must remain inside the weak-limit
interval.  If \(U>1\), it already lies in the compact interval
\([a,1]\Subset(0,U)\).  If \(U=1\), delete \(t_5>1-\delta\).
Ordering forces all five times into \((1-\delta,\tau_n)\),
whose length is at most \(\delta\).  This costs at most
\begin{equation}
\label{nc:upper-endpoint-cost}
 2(C/a)^5\frac{\delta^5}{5!}.
\end{equation}
Choose \(\delta>0\) to make this as small as desired.  No
convergence at the event time is being asserted.

We also remove a thin neighborhood of the boundary of the ordered
time region. Put \(t_0=\tau_n\), and discard paths on which
\(t_5-a<\kappa\) or \(t_{i-1}-t_i<\kappa\) for some
\(1\le i\le5\). Integrating velocities successively with the
bound \(C/a\), their total contribution is at most
\(C_a\kappa\): each of the six time strips has five-dimensional
volume at most \(\kappa\) in the unit cube. Choose \(\kappa>0\)
after the preceding tolerances. On the retained set all consecutive
time gaps are at least \(\kappa\), and the lower time exceeds
\(a+\kappa\). Since \(\tau_n\to1\), the retained times lie in a
fixed compact subset of \(0<t_5<\cdots<t_1<1\) for all large \(n\).

Fix one retained scale/sign tuple. Drop all spatial masks in
constructing its maps, and retain the masks in its nonnegative
integrand.  Each plane increment is bounded, all normals are
bounded above and below, and each cold parameter is bounded.
Apply Lemma~\ref{nc:plane-parameters} successively.  The integral
becomes an integral over five times, fifteen \(L\)-coordinates,
and fifteen cold parameter coordinates.  Its support lies in
a fixed compact subset of
\((0,\infty)^5\times\R^{30}\), and its density is uniformly
bounded: the core amplitudes are at most \(C/a\), the factors
\(|n_i|\zeta(L_i\cdot n_i/|n_i|)\) are bounded, and all retained
ratios \(H(v_i)/H(v_{i-1})\) are bounded.  The masks, time
ordering, and cut conditions are arbitrary measurable factors
between zero and one in this density.

The target \((t_5,x_5,v_5)\) on the actual support lies in a
fixed compact set \(B\) inside the open weak-limit region, away
from the velocity axis.  Set
\(F_n=f_n\ind_B\ind_{\{a\le t\le\tau_n\}}\) and extend it by zero;
it is bounded by two and tends to zero in measure.  The maps converge
in \(C^1\) to the map in Lemma~\ref{lem:five-gain-rank}.
That lemma and Lemma~\ref{nc:uniform-pullback} show that the
integral for this tuple tends to zero.  Notice that the entire
parameter compact need not map into \(B\); only the support of
the measurable density must do so.

There are finitely many tuples, so their sum tends to zero.
The total limsup of the original integral is bounded by the
arbitrarily small errors from the preceding deletions.
Letting those tolerances tend to zero proves the lemma.
\end{proof}

\subsection{Completion of the nonvanishing argument}

\begin{proof}[Completion of the proof of Proposition~\ref{prop:nonzero-limit}]
Assume \(l=0\) and \(g=0\), as in the reduction above.
Choose endpoints \((x_n,v_n)\) at time \(\tau_n\) in the common
full-measure set fixed above such that \(f_n(\tau_n,x_n,v_n)\ge1/2\).
For each fixed \(a\), the residual in
\eqref{nc:five-residual} tends to zero along any subsequence
of these endpoints.  Indeed, Lemma~\ref{nc:endpoint-tightness}
first excludes large velocities and escaping positions.
Otherwise choose convergent endpoints after a periodic
translation.  If the limiting velocity is axial, the first
two gains vanish by Lemma~\ref{nc:two-axial-gains}, and the
last three have bounded mass by \eqref{nc:iteration-norm}.
If it is nonaxial, use \eqref{nc:zero-in-measure} in
Lemma~\ref{nc:five-nonaxial-gains}.  This covers every
subsequence; hence the residual tends to zero for the
original endpoint sequence as well.

Taking the limsup in \eqref{nc:five-residual} at these
endpoints gives
\[
 \frac12
 \le a\sum_{j=0}^4\frac{(C\log(1/a))^j}{j!}.
\]
Finally let \(a\downarrow0\).  The right side tends to zero,
a contradiction.  Throughout this argument \(U\) was allowed
to be infinite, and the case \(U=1\) was handled by
\eqref{nc:upper-endpoint-cost}.  This proves the proposition.
\end{proof}
